\section{Introduction}
\label{sec:intro}

A legal answer is only as good as the authority behind it. An explanation that reads smoothly but
cites the wrong statute, overlooks a controlling precedent, or misses a recent amendment is not a
minor flaw; it is a wrong answer, and one with real consequences for whoever relies on it. This is
what sets legal information systems apart from general-purpose question answering. A useful legal
system must identify the statutes and precedents that actually govern a question, interpret them in
the context that gives them meaning, show how they relate to one another, and present conclusions a
reader can check against the original source. Exact statutory wording matters. A provision's
relationship to its subordinate clauses matters. The authority of a judicial decision matters, and
so do the jurisdiction and date a rule applies to. A response can get any one of these wrong while
still sounding confident and complete, which is precisely what makes legal retrieval harder than
most other retrieval problems. Building a system that respects all of these constraints at once has
occupied legal AI research for more than three decades, and the problem is still open today.

This survey tells that three-decade story as a single, continuous narrative, in the order it
actually happened. Legal artificial intelligence (legal AI) did not arrive through one breakthrough;
it grew through a long sequence of smaller advances, each responding to a specific limitation in the
work before it. The story begins in 1991, when Kevin Ashley's case-based reasoning system HYPO showed
that a computer could compare a new legal dispute with earlier cases and identify the similarities
and differences that make a legal argument \citep{ashley1991}. We take HYPO as our starting point and
move forward year by year, introducing every system, benchmark, and finding in the order the field
actually published it. A reader who wants to understand how today's legal retrieval-augmented
generation (RAG) systems came to exist should be able to retrace the same path the field itself took,
without jumping between decades or topics. So the structure of this survey is the timeline of the
field itself, not a set of themes we arranged for convenience.

Three questions recur across this history, and they give the survey its structure. First,
\emph{how should legal knowledge be represented?} Early systems encoded it as explicit rules and
case comparisons; later systems learned it as patterns in text; the most recent systems treat it as
an external, inspectable corpus that a model retrieves from rather than memorizes. Second, \emph{how
should relevant legal evidence be found?} The field moved from simple keyword matching, to semantic
and hybrid retrieval, to retrieval that tries to preserve the hierarchical structure of a statute or
the reasoning inside a judgment. Third, \emph{how should a system's answer be trusted?} Once
language models became fluent enough to sound authoritative whether or not they were right, the
field had to build ways of checking a generated claim against the evidence that is meant to support
it. These three questions, representation, retrieval, and trust, organize everything that follows.

This survey makes four contributions. First, a \textbf{chronological synthesis} that traces legal
AI's development in the exact order it happened, from case-based reasoning to contemporary
structure-aware retrieval-augmented generation, explaining the problem that motivated each new piece
of work. Second, an \textbf{architectural analysis} that separates document representation,
retrieval, rank fusion, generation, and evidence verification in the order the field introduced
them, and shows how each component feeds into overall reliability. Third, a \textbf{focused
comparative case study} of two systems published in 2026, Legalyze and the data-synthesis approach
of Ren et al., which target different parts of the same pipeline. Fourth, a \textbf{research agenda}
covering open challenges in structure-preserving retrieval, jurisdiction-aware evidence selection,
multilingual legal access, and verifiable answer generation. Table~\ref{tab:rq} summarizes the
research questions behind these four contributions.

\begin{table*}[t]
\centering
\caption{Research questions guiding the survey.}
\label{tab:rq}
\footnotesize
\setlength{\tabcolsep}{5pt}
\renewcommand{\arraystretch}{1.25}
\begin{tabularx}{\textwidth}{>{\raggedright\arraybackslash}X >{\raggedright\arraybackslash}X}
\toprule
\textbf{Research question} & \textbf{Purpose} \\
\midrule
\textbf{RQ1.} How did legal AI evolve, in the order it actually happened, from case-based reasoning to RAG? & Present the principal methodological transitions and the limitations that motivated each one, in sequence. \\
\textbf{RQ2.} How do legal retrieval architectures affect evidence quality? & Explain how systems locate relevant authorities and preserve the context needed to interpret them. \\
\textbf{RQ3.} How should we assess reliability in legal AI? & Distinguish fluent generation from evidence-supported legal answers. \\
\textbf{RQ4.} What challenges remain for jurisdiction-aware and multilingual legal AI? & Identify research opportunities across different legal and linguistic contexts. \\
\bottomrule
\end{tabularx}
\end{table*}

Section~\ref{sec:method} explains how we found, checked, and selected the literature for this
chronological account. The narrative itself then begins in Section~\ref{sec:foundations} with
Ashley's 1991 work on HYPO.
